\changecolours{ScoutGreen}
\section{Taxonomia das Sublinguagens SQL}

% ------------------------------------------------------------------------------
\twocolframe[title={A Linguagem SQL: Visão Global e Estrutura},leftcol=7cm,rightcol=7cm]{%
  \alert{Natureza Declarativa da SQL}\par
  \itemseps[topsepi=1.5mm,itemsepi=2.5mm,labelsep=2mm,leftmargini=4mm]
  \begin{itemize}
    \item \textbf{Linguagem Declarativa:} O desenvolvedor especifica \textbf{o que} deseja obter, e o motor do SGBD define de forma autônoma \textbf{como} navegar pelos blocos físicos de dados.
    \item \textbf{Padrão Internacional:} Normalizada pela ANSI e ISO (ISO/IEC 9075), evoluindo desde o SQL-86 até versões recentes (SQL:2016 e SQL:2023).
    \item \textbf{Subdivisão Funcional:} Embora seja uma linguagem unificada, suas instruções são categorizadas formalmente em cinco sublinguagens especializadas.
  \end{itemize}
}{%
  \alert{As Cinco Famílias de Comandos}\par
  \itemseps[topsepi=1.5mm,itemsepi=2.5mm,labelsep=2mm,leftmargini=4mm]
  \begin{itemize}
    \item \textbf{DDL (\textit{Data Definition}):} Cria, modifica e extingue esquemas e tabelas no Catálogo do Sistema.
    \item \textbf{DML (\textit{Data Manipulation}):} Insere, altera e remove tuplas (linhas) nas tabelas existentes.
    \item \textbf{DQL (\textit{Data Query}):} Realiza buscas, filtros e projeções sem alterar o estado do banco.
    \item \textbf{DCL (\textit{Data Control}):} Administra direitos de acesso, perfis e credenciais de segurança.
    \item \textbf{TCL (\textit{Transaction Control}):} Gerencia a atomicidade e consistência das transações no banco.
  \end{itemize}
}

% ------------------------------------------------------------------------------
\begin{frame}[fragile]{DDL: Linguagem de Definição de Dados}
  \vspace{-0.3cm}
  \begin{columns}[t]
    \begin{column}{0.48\textwidth}
      \alert{Propósito e Impacto no Catálogo}\par
      \vspace{0.05cm}
      \begin{itemize}\setlength{\itemsep}{0pt}\setlength{\parskip}{0pt}
        \item \textbf{Definição:} Modifica o catálogo do SGBD.
        \item \textbf{CREATE:} Instancia tabelas e índices.
        \item \textbf{ALTER:} Modifica estruturas existentes.
        \item \textbf{DROP:} Elimina objetos do banco.
        \item \textbf{TRUNCATE:} Esvazia tabelas rapidamente.
      \end{itemize}
    \end{column}

    \begin{column}{0.48\textwidth}
      \alert{Exemplo Aplicado a Redes}\par
      \vspace{0.05cm}
\begin{lstlisting}[language=SQL,basicstyle=\ttfamily\tiny]
CREATE TABLE interfaces_rede (
  id SERIAL PRIMARY KEY,
  disp_id INT REFERENCES dispositivos(id),
  nome VARCHAR(32) NOT NULL,
  mtu INT DEFAULT 1500
);
ALTER TABLE interfaces_rede 
  ADD COLUMN vlan_id INT;
\end{lstlisting}
    \end{column}
  \end{columns}
\end{frame}

% ------------------------------------------------------------------------------
\begin{frame}[fragile]{DML \& DQL: Manipulação e Consulta}
  \vspace{-0.2cm}
  \begin{columns}[t]
    \begin{column}{0.48\textwidth}
      \alert{DML (Data Manipulation)}\par
      \vspace{0.05cm}
      \begin{itemize}\setlength{\itemsep}{2pt}
        \item \textbf{Finalidade:} Modificar tuplas em tabelas.
        \item \textbf{Comandos:} \texttt{INSERT}, \texttt{UPDATE}, \texttt{DELETE}.
      \end{itemize}
      \vspace{0.05cm}
\begin{lstlisting}[language=SQL,basicstyle=\ttfamily\tiny]
UPDATE interfaces_rede 
SET mtu = 9000 
WHERE nome = 'TenGigE0/1';
\end{lstlisting}
    \end{column}

    \begin{column}{0.48\textwidth}
      \alert{DQL (Data Query)}\par
      \vspace{0.05cm}
      \begin{itemize}\setlength{\itemsep}{2pt}
        \item \textbf{Finalidade:} Extrair dados sem alterar o banco.
        \item \textbf{Comando Canônico:} \texttt{SELECT}.
      \end{itemize}
      \vspace{0.05cm}
\begin{lstlisting}[language=SQL,basicstyle=\ttfamily\tiny]
SELECT d.hostname, i.nome, i.mtu
FROM dispositivos d
JOIN interfaces_rede i 
  ON i.dispositivo_id = d.id
WHERE i.mtu = 9000;
\end{lstlisting}
    \end{column}
  \end{columns}
\end{frame}

% ------------------------------------------------------------------------------
\begin{frame}[fragile]{DCL \& TCL: Segurança e Transações}
  \vspace{-0.3cm}
  \begin{columns}[t]
    \begin{column}{0.48\textwidth}
      \alert{DCL (Data Control)}\par
      \vspace{0.05cm}
      \begin{itemize}\setlength{\itemsep}{2pt}
        \item \textbf{Finalidade:} Gerenciar privilégios e perfis.
        \item \textbf{Comandos:} \texttt{GRANT}, \texttt{REVOKE}.
      \end{itemize}
      \vspace{0.05cm}
\begin{lstlisting}[language=SQL,basicstyle=\ttfamily\tiny]
GRANT SELECT ON interfaces_rede 
  TO operador_noc;
REVOKE DROP ON ALL TABLES 
  FROM operador_estagiario;
\end{lstlisting}
    \end{column}

    \begin{column}{0.48\textwidth}
      \alert{TCL (Transaction Control)}\par
      \vspace{0.05cm}
      \begin{itemize}\setlength{\itemsep}{2pt}
        \item \textbf{Finalidade:} Assegurar propriedades ACID.
        \item \textbf{Comandos:} \texttt{COMMIT}, \texttt{ROLLBACK}, \texttt{SAVEPOINT}.
      \end{itemize}
      \vspace{0.05cm}
\begin{lstlisting}[language=SQL,basicstyle=\ttfamily\tiny]
SAVEPOINT sp_ponto_seguro;
-- atualizacao de teste
ROLLBACK TO sp_ponto_seguro;
COMMIT;
\end{lstlisting}
    \end{column}
  \end{columns}
\end{frame}
